\documentclass[10pt,a4paper]{amsart}
\usepackage{amsmath,amssymb,amsthm,graphicx}
% Retained minimal preamble for the Stacks Project chapter more-algebra.tex (More on Algebra).
% The chapter file pinned-original-source.tex carries only \input{preamble} on its first line, so no
% preamble text lives inside the chapter itself.  The Stacks Project uses one single shared preamble for
% every chapter; the snapshot stages it identically as sources/stacks-chow/preamble.tex and
% sources/stacks-etale-cohomology/preamble.tex.  The block below is copied mechanically, line for line,
% from sources/stacks-chow/preamble.tex lines 14-18 (the Xy-pic setup) and 154-217 (the native theorem
% declarations together with the mathematical macros).  Nothing was added, removed or redefined; the
% wrapper class and the amsmath/amssymb/amsthm/graphicx base packages are supplied by the assembler.
% Chapter: more-algebra.tex -- More on Algebra.
% For commutative diagrams we use Xy-pic
\usepackage[all]{xy}

% We use 2cell for 2-commutative diagrams.
\xyoption{2cell}

% Theorem environments.
%
\theoremstyle{plain}
\newtheorem{theorem}[subsection]{Theorem}
\newtheorem{proposition}[subsection]{Proposition}
\newtheorem{lemma}[subsection]{Lemma}

\theoremstyle{definition}
\newtheorem{definition}[subsection]{Definition}
\newtheorem{example}[subsection]{Example}
\newtheorem{exercise}[subsection]{Exercise}
\newtheorem{situation}[subsection]{Situation}

\theoremstyle{remark}
\newtheorem{remark}[subsection]{Remark}
\newtheorem{remarks}[subsection]{Remarks}

\numberwithin{equation}{subsection}

% Macros
%
\def\lim{\mathop{\mathrm{lim}}\nolimits}
\def\colim{\mathop{\mathrm{colim}}\nolimits}
\def\Spec{\mathop{\mathrm{Spec}}}
\def\Hom{\mathop{\mathrm{Hom}}\nolimits}
\def\Ext{\mathop{\mathrm{Ext}}\nolimits}
\def\SheafHom{\mathop{\mathcal{H}\!\mathit{om}}\nolimits}
\def\SheafExt{\mathop{\mathcal{E}\!\mathit{xt}}\nolimits}
\def\Sch{\mathit{Sch}}
\def\Mor{\mathop{\mathrm{Mor}}\nolimits}
\def\Ob{\mathop{\mathrm{Ob}}\nolimits}
\def\Sh{\mathop{\mathit{Sh}}\nolimits}
\def\NL{\mathop{N\!L}\nolimits}
\def\CH{\mathop{\mathrm{CH}}\nolimits}
\def\proetale{{pro\text{-}\acute{e}tale}}
\def\etale{{\acute{e}tale}}
\def\QCoh{\mathit{QCoh}}
\def\Ker{\mathop{\mathrm{Ker}}}
\def\Im{\mathop{\mathrm{Im}}}
\def\Coker{\mathop{\mathrm{Coker}}}
\def\Coim{\mathop{\mathrm{Coim}}}

% Boxtimes
%
\DeclareMathSymbol{\boxtimes}{\mathbin}{AMSa}{"02}

%
% Macros for moduli stacks/spaces
%
\def\QCohstack{\mathcal{QC}\!\mathit{oh}}
\def\Cohstack{\mathcal{C}\!\mathit{oh}}
\def\Spacesstack{\mathcal{S}\!\mathit{paces}}
\def\Quotfunctor{\mathrm{Quot}}
\def\Hilbfunctor{\mathrm{Hilb}}
\def\Curvesstack{\mathcal{C}\!\mathit{urves}}
\def\Polarizedstack{\mathcal{P}\!\mathit{olarized}}
\def\Complexesstack{\mathcal{C}\!\mathit{omplexes}}
% \Pic is the operator that assigns to X its picard group, usage \Pic(X)
% \Picardstack_{X/B} denotes the Picard stack of X over B
% \Picardfunctor_{X/B} denotes the Picard functor of X over B
\def\Pic{\mathop{\mathrm{Pic}}\nolimits}
\def\Picardstack{\mathcal{P}\!\mathit{ic}}
\def\Picardfunctor{\mathrm{Pic}}
\def\Deformationcategory{\mathcal{D}\!\mathit{ef}}

% --- appended 2026-10-12 (Chief): the project own definitions of the reference/slogan/history
% display environments, copied mechanically from the shared Stacks preamble lines 8-12
% (\usepackage{verbatim} plus the three \newenvironment lines, which the chapter preamble keeps
% verbatim from the project preamble). Nothing added, removed or redefined: these environments
% carry no mathematical content and the project itself typesets them as comments. Without these
% three lines a counted statement whose own text uses them fails with Environment slogan undefined.
\usepackage{verbatim}
\newenvironment{reference}{\comment}{\endcomment}
\newenvironment{slogan}{\comment}{\endcomment}
\newenvironment{history}{\comment}{\endcomment}
\usepackage{cleveref}
\crefname{originallemma}{original lemma}{original lemmas}
\Crefname{originallemma}{Original lemma}{Original lemmas}
\makeatletter
\newlabel{lemma-pull-pseudo-coherent}{{0650}{source}{}{sourcepointer-1}{}}
\newlabel{lemma-pull-pseudo-coherent@cref}{{[originallemma][0][]0650}{[1][0][]source}{}{}{}}
\makeatother
\crefname{originallemma}{original lemma}{original lemmas}
\Crefname{originallemma}{Original lemma}{Original lemmas}
\makeatletter
\newlabel{lemma-pseudo-coherent}{{064U}{source}{}{sourcepointer-2}{}}
\newlabel{lemma-pseudo-coherent@cref}{{[originallemma][0][]064U}{[1][0][]source}{}{}{}}
\makeatother
\crefname{originallemma}{original lemma}{original lemmas}
\Crefname{originallemma}{Original lemma}{Original lemmas}
\makeatletter
\newlabel{lemma-lift-complex-projectives}{{09AR}{source}{}{sourcepointer-3}{}}
\newlabel{lemma-lift-complex-projectives@cref}{{[originallemma][0][]09AR}{[1][0][]source}{}{}{}}
\makeatother
\title{Pseudo-coherence descends along a surjection with nilpotent kernel}
\author{The Stacks Project Authors}
\date{stacks-project:more-algebra.tex:local-snapshot-20261012; GNU FDL 1.2 or later}
\begin{document}
\maketitle
\setcounter{subsection}{1}
\noindent\textit{Editorial note.} All delivered mathematical text is the author's own, copied from the pinned Stacks Project chapter file more-algebra.tex; no statement or proof step was rewritten, reordered or removed. The counted span is the author environment and its own proof environment. References inside the span are handed to the assembler as navigation stubs that carry the canonical tag of the cited tag only; whether each of them is load-bearing is recorded on the stub itself.\par\medskip
\noindent\textit{Allowed standard prerequisites (editorial).}\begin{itemize}
\item A pseudo-coherent complex can be represented by a bounded above complex of finite free modules. invoked at proof lines 19505-19507. Produces the finite free complex E whose terms are the reductions of the complex to be built on the other side.
\item A bounded above complex over the quotient with a lift of its terms can be lifted to a bounded above complex of projective modules with the prescribed reductions. invoked at proof lines 19507-19511. Supplies the complex P of projective modules lifting E, which is the candidate representing object on the other side.
\item Nakayama's lemma lets one conclude that a projective module over the ring whose reduction to the quotient is finite free is itself finite free. invoked at proof lines 19512-19513. Upgrades the projectivity of the terms of P to finite freeness, which is what pseudo-coherence of the lifted object requires.
\end{itemize}\medskip
\ifdefined\hypertarget\hypertarget{sourcepointer-1}{}\fi
\noindent\textit{Source reference.} Original lemma 0650: sources/stacks-more-algebra/pinned-original-source.tex, source lines 16345--16345. Used at proof line 19504 for the direction that pseudo-coherence of the lifted object implies pseudo-coherence of K; the converse direction, which is the substance, is inside the counted span.\par\medskip
\ifdefined\hypertarget\hypertarget{sourcepointer-2}{}\fi
\noindent\textit{Source reference.} Original lemma 064U: sources/stacks-more-algebra/pinned-original-source.tex, source lines 16048--16048. Invoked at proof lines 19505-19507 to represent K by a finite free complex; the lifting and the Nakayama argument built on that representation are inside the counted span.\par\medskip
\ifdefined\hypertarget\hypertarget{sourcepointer-3}{}\fi
\noindent\textit{Source reference.} Original lemma 09AR: sources/stacks-more-algebra/pinned-original-source.tex, source lines 19469--19469. Used at proof lines 19507-19511 to lift the finite free complex to a projective one refining the given reduction; checking that the lifted complex has the right reductions and the final Nakayama step are inside the counted span.\par\medskip
\begin{lemma}
\label{lemma-check-pseudo-coherent-modulo-nilpotent}
Let $R' \to R$ be a surjective ring map whose kernel is a nilpotent ideal.
Let $K' \in D(R')$ and set $K = K' \otimes_{R'}^\mathbf{L} R$.
Then $K$ is pseudo-coherent if and only if $K'$ is pseudo-coherent.
\end{lemma}

\begin{proof}
One direction follows from Lemma \ref{lemma-pull-pseudo-coherent}.
For the other direction, assume $K$ is pseudo-coherent. Then by
Lemma \ref{lemma-pseudo-coherent} we can represent $K$ by a bounded
above complex $E^\bullet$ of finite free $R$-modules. By
Lemma \ref{lemma-lift-complex-projectives} we can represent
$K'$ by a bounded above complex $P^\bullet$ of projective
$R'$-modules such that $P^n \otimes_{R'} R = E^n$.
By Nakayama's lemma we see that $P^n$ is finite free
and we conclude that $K'$ is pseudo-coherent as well.
\end{proof}

\end{document}
